\chapter{The Arrow of Time}
\label{chap:timearrow}

\section{The problem}

Standard physics is time-symmetric at the microscopic level. CPT is a theorem;
the equations of motion (Schr\"odinger, Dirac, Einstein) are reversible. Yet we
experience a clear forward arrow of time: eggs break and do not unbreak, heat
flows from hot to cold, memories accumulate of the past but not the future.

Standard cosmology explains this through the \emph{Past Hypothesis}: the universe
began in an extraordinarily low-entropy state, and the arrow of time is the
gradient pointing away from that initial state. But this is uncomfortable
because:

\begin{itemize}
  \item the low-entropy initial state is postulated, not explained;
  \item Boltzmann's original suggestion (we are a rare fluctuation in eternal
    equilibrium) does not work---it predicts that Boltzmann brains should
    dominate over ordered observers, contradicting observation.
\end{itemize}

\section{\SCT resolution}

In \SCT, the arrow of time is automatic:

\begin{enumerate}
  \item \textbf{The Cartasis membrane is a low-entropy boundary.} At the bounce,
    matter is highly compressed ($\sim 10^{50}\,\mathrm{kg\,m^{-3}}$) and highly
    ordered (specific chirality bias inherited from the parent). Gravitational
    entropy is low because matter is uniformly compressed rather than clumped
    into black holes.
  \item \textbf{The post-bounce evolution is toward higher entropy.} Expansion,
    cooling, structure formation, stellar evolution, eventual black-hole
    formation, eventual Hawking evaporation---all standard cosmological evolution
    proceeds toward higher entropy.
  \item \textbf{The arrow of time is the gradient between these.} Forward time
    $=$ away from bounce $=$ increasing entropy. Backward time $=$ toward bounce
    $=$ decreasing entropy.
\end{enumerate}

The Past Hypothesis is no longer a postulate; it is a \emph{consequence} of the
bounce being a low-entropy boundary.

\section{Why doesn't this just push the problem back?}
\label{sec:forced-low-entropy}

One might object: ``fine, but now you need to explain why the bounce is
low-entropy.'' The answer is that the bounce \emph{cannot} be high entropy---it
is the geometric configuration where matter is maximally compressed within
Cartasis density limits. At Cartasis density, you have already gone beyond
conditions where complex structure (which carries information and entropy) can
exist. The bounce is necessarily a state of high concentration, low complexity,
and low information content---which means low gravitational entropy.

In other words, the bounce membrane is mechanically forced to be low entropy by
the physics of Einstein--Cartan at extreme density. You cannot have a
high-entropy bounce because high entropy requires complex structure, and complex
structure cannot exist at Cartasis density. This converts the Past Hypothesis
from ``an unexplained initial condition'' to ``a necessary consequence of bounce
physics.''

\section{Global vs.\ local time}

In \SCT, time has a particular structure. There is a \emph{local arrow within
each universe}: each universe has its own internal time direction, set by its
bounce-to-future thermodynamic gradient, and observers within a universe
experience their universe's arrow. There is \emph{no global supraverse time
arrow}: different universes have arrows pointing in different directions in
supraverse coordinates, and the supraverse as a whole is time-symmetric. And
there is a \emph{static block universe at the supraverse level}: from the
supraverse perspective, all of supraverse spacetime exists as a 4D block;
``evolution'' is a feature of being inside a specific universe, not a feature of
the supraverse itself. This is the standard relativistic block-universe view,
extended to the supraverse.

\section{Does the supraverse need a time arrow? No.}
\label{sec:no-supraverse-arrow}

It is worth stating the strongest form of this, because it is a place where the
framework is \emph{more} parsimonious than it first appears, not less. An arrow of
time is never present in the microscopic laws---QFT, general relativity, and
Einstein--Cartan are all CPT-symmetric, time-reversal invariant up to the CP phase.
An arrow is not a law but a \emph{boundary condition}: it appears only where a system
is given a low-entropy edge. The supraverse, taken as a whole, has no such edge, and
\emph{needs none}. Positing a global supraverse arrow would add unexplained
structure---a cosmic ``now'' with no physics behind it; declining to posit one costs
nothing and is the honest, minimal choice.

So the right question about the supraverse is not ``which way does its time run?''
but the timeless, thermodynamic one: \emph{what is the stable state of the void?} That
is an equilibrium question, an eigenstate question---``time is what clocks measure,''
and an equilibrium void has no clocks. The answer (Chapter~\ref{chap:supraverse}) is
the condensed foam: the stationary, maximum-entropy configuration of an infinite
self-gravitating quantum vacuum. No arrow enters anywhere in posing or answering it.

\paragraph{Where every arrow comes from, then.} Each bounce is a low-entropy
boundary, and a thermodynamic arrow is manufactured \emph{there}---pointing away from
the bounce, toward higher entropy, independently at every membrane. Time is therefore
\emph{emergent, local, and created per universe}, not inherited from any
supraverse-level clock. There is no contradiction between a timeless void and the
manifest passage of time we experience: we experience time precisely because we live
\emph{inside} one of the void's fluctuation-condensations, on the high-entropy side of
its bounce.

\paragraph{The apparent paradox of ``process'' in a timeless void.} One might object
that the supraverse \emph{does} seem to require time: universes ``nucleate,''
equilibrium is ``established,'' the foam ``coarsens.'' But these are descriptions from
\emph{within} the internal time of a condensation, projected onto the whole. The void
does not \emph{evolve toward} equilibrium in some external time; the equilibrium
\emph{is} the timeless stationary state, and ``nucleation'' is simply the statement
that the equilibrium ensemble \emph{contains} fluctuation-condensations---the foam---as
part of what the stationary state \emph{is}. No global clock ticks. The language of
process is an unavoidable artifact of describing a timeless structure in the
time-laden words available to observers who are themselves internal to it. This
dissolves the apparent tension (\Cref{app:openq}, Q29) rather than leaving it open: the
block-universe void and the ``dynamical-sounding'' foam are the same thing, seen from
outside time and from inside it.

\section{Why we experience forward time}
\label{sec:why-forward}

Observers exist within specific universes. Observers require metabolism
(irreversible chemistry, time-asymmetric), memory (information accumulation,
requiring an entropy gradient), and computation/thought (irreversible logic
gates, dissipative). All of these require being far from a low-entropy boundary
and moving away from it. Observers can only exist where there is a thermodynamic
gradient to extract work from.

So observers necessarily exist in regions with a clear arrow of time pointing
away from a low-entropy boundary (the bounce). \textbf{We experience forward time
because we could not exist otherwise.} It is not a metaphysical mystery; it is an
anthropic-style consequence of where observers can exist within the supraverse
structure.

\section{Why we are not Boltzmann brains}
\label{sec:not-brains}

The forward-time argument has a sharp corollary that disposes of the oldest
objection to fluctuation cosmologies---and, with it, of the suspicion that we are
a simulation. The objection (chapter opening): a mind with memories is a
low-entropy ordered state, so nucleating one directly out of the vacuum looks far
cheaper than growing an entire universe, and a \emph{typical} observer should
then be such a fluctuation---a ``Boltzmann brain''---rather than an evolved
inhabitant of a real cosmos. Cast in the de Sitter Boltzmann weights of
Chapter~\ref{chap:supraverse}, the bare nucleation costs seem to confirm it:
\begin{equation}
  \underbrace{10^{-10^{69}}}_{\text{a brain }(\sim 1\,\mathrm{kg})}
  \;\gg\;
  \underbrace{10^{-10^{84}}}_{\text{an OG seed }(\sim 10^{15}\,\mathrm{kg})}
  \;\gg\;
  \underbrace{10^{-10^{123}}}_{\substack{\text{a finished low-entropy}\\\text{Big Bang \citep{penrose2010}}}},
  \label{eq:three-towers}
\end{equation}
the brain winning by the unbridgeable factor $10^{10^{84}}$. And no rescue
helps: not the head-count of real observers a universe produces, not the
recursive tree of descendant universes, not the longer lifetime of a real
observer over a brain's single flickering thought. Each of those is a
\emph{single} power tower---a number of at most a few thousand digits---and none
can dent a \emph{double} tower of height $10^{84}$.

The resolution is not to win the count; it is to see that it is the wrong count,
because \emph{the brain and the bare Big Bang cannot arise from the primordial
void at all.} A Boltzmann brain---or, equivalently, the optimal computer running
a simulation of one---is an information-processing system, and
Section~\ref{sec:why-forward} already fixed what such a system requires:
metabolism, memory, and computation are dissipative, and dissipation needs a
thermodynamic arrow and a low-entropy reservoir to discharge into. The void
supplies neither. It is the maximum-entropy de Sitter \emph{eigenstate}
(Chapter~\ref{chap:supraverse})---stationary, arrowless, with no dynamics to
fluctuate and no gradient to compute against. A brain is parasitic on an arrow it
cannot itself create; the $10^{-10^{69}}$ was computed as though the void were a
fluctuating thermal bath, which it is not~\citep{boddy2014}.

What the eigenstate \emph{can} do is decay. Nucleation from a stationary state is
a genuine quantum transition---a gravitational instanton, the object whose action
Chapter~\ref{chap:supraverse} computes---and it yields only the \emph{generic}
state of the mass-energy it creates: a structureless, maximal-entropy,
black-hole-like blob. Such a blob is exactly a viable OG seed. It carries no
improbable order of its own; yet when it bounces, the bounce is \emph{forced} to
be low-entropy (Section~\ref{sec:forced-low-entropy}), manufacturing order and
arrow for free, downstream of the tunnelling. A brain is not a blob. It is a
low-entropy \emph{organized} state, and nucleation cannot deliver
organization---only generic high-entropy lumps---while the arrowless void cannot
assemble it afterward. So the void makes seeds, never brains. The seed's
nucleation rate carries a de Sitter Boltzmann factor $e^{-I}$, but the underlying
event is a tunnelling \emph{decay}, not a within-vacuum thermal
fluctuation---which is precisely why it fires from the eigenstate when a brain
cannot.

The same blade cuts Penrose's bare Big Bang, and this is the deep part. A hot
Big Bang sitting in its $1$-in-$10^{10^{123}}$ low-entropy starting
state~\citep{penrose2010} is \emph{also} an ordered, arrow-bearing
configuration---posited with no mechanism to produce it. It presupposes exactly
the low-entropy past it is meant to explain, so it too cannot be nucleated from
the void; it can only be put in by hand. A cosmology whose sole origin story is
``the initial state happened to be that special'' is therefore, structurally, the
simulation hypothesis: it declares us an un-generated fluke of order with no
maker. \emph{If the bare Big Bang were the only available beginning, we would
indeed be a simulation.} \SCT\ escapes because it never posits the order: it pays
$\sim 10^{-10^{84}}$ to nucleate a generic seed and lets the bounce \emph{force}
the $10^{123}$ worth of low-entropy order at no extra cost---converting Penrose's
improbability-of-order from an unexplained initial condition into a consequence
of Einstein--Cartan geometry at the membrane.

Two levels then close the problem. At the \emph{foundation}, nothing mind-like
can bottom out the stack; only an arrow-creating bounce can, so the base layer of
reality is real universes, not a sea of fluctuations. \emph{Inside} a universe
the arrow exists, so brains \emph{can} fluctuate---but only over the de Sitter
recurrence time $\sim 10^{10^{69}}$ years, and \SCT's dark energy is finite
(parent accretion, Chapter~\ref{chap:darksector}): it switches off when the
parent finishes evaporating, $\sim 10^{136}$ years, recycling the matter into new
bounces long before a single interior brain could form. A $\Lambda$CDM model with
a truly eternal cosmological constant has no such escape---its endless de Sitter
future is a literal Boltzmann-brain factory---so this is a respect in which \SCT\
is \emph{strictly better behaved} than the standard model, not merely different.

\paragraph{Epistemic status.} The de Sitter measure problem---whether and how a
stationary vacuum ``fluctuates''---is not settled, and a fully rigorous \SCT\
measure does not yet exist (\Cref{app:openq}). What is robust is the qualitative
asymmetry: nucleation yields generic blobs, blobs bounce into arrowed universes,
and organized low-entropy minds presuppose an arrow they cannot source from an
arrowless eigenstate. The cognitive-stability argument seals it from the other
side---any theory predicting that typical observers are fluctuations dissolves
its own evidential basis, since a fluctuated mind has no warrant to trust the
memories, records, and reasoning that built the theory, including the three
numbers in \eqref{eq:three-towers}. We are entitled to conclude we are real not
by counting, but because the alternative is self-refuting and because the void's
only product is the real thing. Reproducible as
\texttt{sims/src/cartasis\_sims/boltzmann\_brains.py}.

\section{Implications for cosmology}

\subsection{No initial-conditions problem}

Standard cosmology has to postulate that the early universe was low-entropy and
homogeneous. \SCT derives both:

\begin{itemize}
  \item low entropy: forced by bounce physics;
  \item homogeneity: bounce dynamics smooth out parent inhomogeneities, plus
    rapid expansion redshifts away small-scale perturbations (similar to
    inflation's mechanism).
\end{itemize}

\subsection{Eventual heat death is local}

In standard cosmology, the universe heads toward heat death (maximum entropy, no
usable energy). In \SCT, this happens too, locally---our universe will eventually
reach a state of maximum entropy where structure no longer forms. But the
\emph{supraverse} does not approach heat death. The supraverse is already at its
equilibrium configuration. Our universe's eventual heat death is just one bubble
reaching its local equilibrium; the supraverse as a whole maintains its
statistical distribution.

\subsection{Universes evaporating back to parents}

As discussed in Chapter~\ref{chap:baby}, the parent's Hawking evaporation
eventually shrinks the Cartasis membrane and the baby's connection to the parent
fades. At some point the baby's matter content (now entangled with Hawking
radiation in the parent's exterior) is effectively transferred back to the
supraverse vacuum. This is the \emph{cosmological} analog of an individual
universe reaching heat death: not local equilibrium but reabsorption into the
supraverse substrate. The arrow of time within the baby ends when the baby ends.

\section{Connection to gravitational entropy}

Penrose's longstanding emphasis on gravitational entropy is central
here~\citep{penrose1979}. The key insights:

\begin{enumerate}
  \item \textbf{Uniform vacuum is low gravitational entropy.} Smooth metric, no
    horizons, no information stored geometrically.
  \item \textbf{Black holes are maximum entropy.} Bekenstein--Hawking
    $S \propto A/4\ell_P^2$. A black hole's entropy can exceed the entropy of all
    the matter that fell into it.
  \item \textbf{The early universe was low gravitational entropy} despite being
    hot---because gravity was smooth, with no horizons and no clumping. Penrose's
    Weyl curvature hypothesis: the universe started with vanishing Weyl
    curvature.
  \item \textbf{The universe evolves toward higher gravitational entropy} as
    structure forms, especially as black holes form.
\end{enumerate}

In \SCT, the bounce is necessarily a low-Weyl-curvature configuration (matter is
uniformly at Cartasis density, no clumps). This automatically gives the
low-gravitational-entropy initial condition that Penrose's hypothesis requires.
The arrow of time follows.

\section{Connection to gravitational decoherence}

\emph{Gravitational decoherence} is a mechanism in which large mass--energy
superpositions decohere into definite states via gravitational
coupling~\citep{diosi1989,penrose1996}. This is a kind of ``miniature
universe-splitting''---the different branches become causally weakly-coupled
regions of the gravitational field. In \SCT, gravitational decoherence is the
small-scale version of the same physics as bounce-driven universe production.
Both are mechanisms where gravity separates different mass--energy configurations
into causally distinct regions. The full continuum is:

\begin{itemize}
  \item tiny mass--energy differences: weak gravitational decoherence, branches
    remain quantum-mechanically connected;
  \item moderate: stronger decoherence, branches become effectively separate;
  \item extreme (Cartasis density): bounce, branches become full baby universes.
\end{itemize}

The arrow of time within each branch, each universe, is set by the same
mechanism: gravity carves out a low-entropy starting configuration, and time
flows away from it.

\section{The arrow and matter are CPT-locked}

The thermodynamic arrow above is set by the low-entropy bounce boundary and would
exist even in a perfectly matter/antimatter-symmetric universe; it is \emph{not}
caused by the baryon asymmetry. But the two are co-sourced at the same membrane
and, at the supraverse level, locked together. The chiral-vortical filter that
makes the matter excess (Chapter~\ref{chap:chirality}) is CP-violating, and by
the CPT theorem CP violation is T violation: matter genesis is intrinsically
time-asymmetric. More tellingly, the void is CPT-symmetric, so for every
universe that is matter-biased with a forward thermodynamic arrow there is a
CPT-image partner that is antimatter-biased with the opposite arrow---and the two
are the \emph{same} object viewed through CPT. The symmetric quantity is the
product of matter-sign and arrow-sign, and observership selects the diagonal:
every observer necessarily finds ``matter $+$ forward time'' (their mirror's
inhabitants say the same of their ``antimatter $+$ backward'' world). Matter and
time travel in pairs not because one causes the other, but because the CPT
structure of the void binds their signs.
